\scanpage{321}

\unsure{Un} élément de $\operatorname{Discr\acute{e}t}^{\unsure{\prime}}(\pi)$ s'appelle une
classe de discrétifications de $\pi$\inserted{\unsure{ou} \unsure{pré-discrétification} de $\pi$}.
Si $\pi_0$ est une discrétification, on désigne sa classe par
$\pi_0^{\unsure{\natural}}$. L'ens. des classes de discrétifications sur $\pi$
est un \unsure{espace} hom. $\widehat{\widehat{\unsure{\mathfrak{S}|\mathcal{G}}}}(\pi)$, le stabilisateur
de $\pi_0^{\unsure{\natural}}$ \unsure{s'identifie à|s'identifiant} : l'adhérence de
$\unsure{\mathfrak{S}|\mathcal{G}}(\pi_0)=\operatorname{Aut}_{\mathrm{lac}}(\pi_0)$ dans
$\widehat{\widehat{\unsure{\mathfrak{S}|\mathcal{G}}}}(\pi)$, ou encore :
$\widehat{\unsure{\mathfrak{S}|\mathcal{G}}}(\pi_0)$.
Celle-ci contient toujours $\pi$.
\[
\operatorname{Aut}_{\mathrm{lac}}(\pi,\pi_0^{\unsure{\natural}})
\simeq\widehat{\unsure{\mathfrak{S}|\mathcal{G}}}(\pi_0)
\]

\deleted{On a} pour tte discrétification $\pi_0^{\unsure{\natural}}$, le sous-groupe de
$\widehat{\widehat{\unsure{\mathfrak{S}|\mathcal{G}}}}(\pi)$ des \unsure{autom.} à lacets qui fixent
$\pi_0^{\unsure{\natural}}$ est noté $\widehat{\unsure{\mathfrak{S}|\mathcal{G}}}(\pi_0^{\unsure{\natural}})$,
c'est donc l'image inverse d'un ss-groupe de
$\widehat{\widehat{\mathcal{T}}}(\pi)$, noté
$\widehat{\mathcal{T}}(\pi_0^{\unsure{\natural}})$. On a donc une inclusion de
structures d'extensions

(10)
\diagram{\begin{tikzcd}[ampersand replacement=\&]
1 \arrow[r] \& \pi \arrow[r] \&
\widehat{\unsure{\mathfrak{S}|\mathcal{G}}}(\pi_0^{\unsure{\natural}}) \arrow[r] \arrow[d,hook] \&
\widehat{\mathcal{T}}(\pi_0^{\unsure{\natural}}) \arrow[r] \arrow[d,hook] \& 1 \\
1 \arrow[r] \& \pi \arrow[r] \&
\widehat{\widehat{\unsure{\mathfrak{S}|\mathcal{G}}}}(\pi) \arrow[r] \&
\widehat{\widehat{\mathcal{T}}}(\pi) \arrow[r] \& 1
\end{tikzcd}}

\margin{\uline{NB} L'ens des classes de discrétifications de $\pi$ se décrit
\illegible{2} du groupe \uline{ext} $\pi$, comme
$\operatorname{Isom\,ext}(\widehat{\Pi}_{g,\nu},\pi)$\deleted{\illegible{1}}
divisé par $\widehat{\mathcal{T}}_{g,\nu}$.
On a d'ailleurs \unsure{une flèche} de degré $2$
\diagram{\begin{tikzcd}[ampersand replacement=\&]
\operatorname{Isom\,ext}(\widehat{\Pi}_{g,\nu},\pi)/\widehat{\mathcal{T}}_{g,\nu}^{+}
\arrow[d] \\
\operatorname{Isom\,ext}(\widehat{\Pi}_{g,\nu},\pi)/\widehat{\mathcal{T}}_{g,\nu}
\end{tikzcd}}
ce qui permet pour \unsure{tte|les} \unsure{clas.} de discrétifications, de définir
ses deux \uline{\unsure{orientations}}, et parler des
discrétifications \uline{orientées}.}

\mspage{320}

\scanpage{322}
\mspage{163}

(On montre par voie arithmético-géométrique que l'inclusion
$\widehat{\mathcal{T}}\longrightarrow\widehat{\widehat{\mathcal{T}}}$
n'est jamais un iso.) On trouve ainsi une application
\[
\tag{11}
\operatorname{Discr\acute{e}t}^{\unsure{\prime}}(\pi)\longrightarrow
\text{Ss-groupes fermés de }\widehat{\widehat{\mathcal{T}}}(\pi)
\]
évidemment compatible à l'action de
\deleted{$\unsure{\widehat{\widehat{\mathcal{T}}}}$}
$\widehat{\widehat{\unsure{\mathfrak{S}|\mathcal{G}}}}(\pi)$ (transport de structure) -- opérant
à droite via $\widehat{\widehat{\mathcal{T}}}(\pi)$ et ses autom. int. sur lui-même.

On trouve ainsi une classe de conjugaison, bien déterminée \unsure{de|des}
ss-groupes fermés de $\widehat{\widehat{\mathcal{T}}}(\pi)$,
qu'on appelle ses sous-groupes de Teichmüller ``\uline{géométriques}''.
\deleted{Il n'est nullement improbable d'ailleurs que
$\widehat{\mathcal{T}}(\pi_0^{\unsure{\natural}})$ soit \illegible{1}}
Il se pourrait d'ailleurs que $\widehat{\mathcal{T}}_{g,\nu}$ soit
son propre normalisateur dans $\widehat{\widehat{\mathcal{T}}}_{g,\nu}$
-- équivalent à l'assertion que (11) est \unsure{injective|bijective|injection} :
La donnée d'une classe de discrétifications\inserted{de $\pi$}
serait équivalente à celle d'un ss-groupe de Teichmüller géométrique dans
\deleted{$\unsure{\widehat{\widehat{\mathcal{T}}}}$}
son groupe de Teichmüller \unsure{étendu} $\widehat{\widehat{\mathcal{T}}}$.

\margin{c'est complètement déconnant et ultra-faux
-- un moment d'égarement ! Cela apparaît clairement par la suite\ldots
La question \unsure{judicieuse} (avec laquelle \unsure{j'ai} dû
\unsure{sur le coup} confondre) \unsure{c'est} si $\widehat{\mathcal{T}}$ est
\uline{invariant} dans $\widehat{\widehat{\mathcal{T}}}$, i.e. si
le normalisateur est $\widehat{\widehat{\mathcal{T}}}$ \unsure{tt|tout} \unsure{entier} !}

\mspage{321}

\scanpage{323}

Notons qu'on a les hom canoniques
\[
\tag{12}
\widehat{\widehat{\mathcal{T}}}\deleted{\unsure{(\pi_0^{\natural})}}
\xrightarrow{\chi}\widehat{\unsure{\mathbb{Z}}}^{*}
\]
\[
\tag{13}
\widehat{\widehat{\mathcal{T}}}\longrightarrow
\unsure{\mathfrak{S}_{I}|\mathcal{G}_{I}}
\]
\margin{où $I=I(\pi)$ est \unsure{l'ens|l'ens.} des classes de conj. de ss-groupes à lacets de $\pi$}
induisant sur $\widehat{\mathcal{T}}(\pi_0^{\unsure{\natural}})$ des
\deleted{\unsure{continus}} hom correspondants -- d'où la définition des
ss groupes $\widehat{\widehat{\mathcal{T}}}^{!}$,
$\widehat{\widehat{\mathcal{T}}}^{+}$,
$\widehat{\widehat{\mathcal{T}}}^{!+}$ et de $\overline{\mathrm m}$ pour
$\widehat{\mathcal{T}}$. Noter que $\chi$ ne prend sur $\widehat{\mathcal{T}}$
que les valeurs $\pm1$.

Le sous-groupe $\widehat{\mathcal{T}}^{+}$ des autom. ext. des
\unsure{grps|groupes} à lacets profinis $\pi$,
\deleted{\illegible{2}}\inserted{fixant} une classe de discrétifications
$\pi_0^{\unsure{\natural}}$ et de multiplicateur $+1$, joue un rôle très particulier.
À l'opposé de ce \unsure{qu'on} peut conjecturer sur $\widehat{\mathcal{T}}$
(dont $\widehat{\mathcal{T}}^{+}$ est un sous-groupe \unsure{ouvert|invariant} d'indice $2$),
on supposerait plutôt que $\widehat{\mathcal{T}}^{+}$ est invariant dans
$\widehat{\widehat{\mathcal{T}}}$. En tt état de cause,
les ss-groupes ainsi définis dans $\widehat{\widehat{\mathcal{T}}}$
\deleted{qui}\inserted{via} \unsure{classes} de discrétifications de $\pi$
\deleted{\illegible{1}} (peut-être

\mspage{322}

\scanpage{324}
\mspage{164}

n'y en a-t-il qu'un seul et unique!) s'appelleront les sous-groupes de Teichmüller géométriques \uline{stricts}. Le choix d'un tel sous-groupe \inserted{de $\widehat{\widehat{\mathcal{T}}}$} est, en tt état de cause, un élément de structure \unsure{nettement} plus faible que celle d'une classe de \unsure{discrétifications}, et $\overline{\mathrm m}$ que celle d'un ss-groupe de Teichmüller géométrique (pas strict).

Pour préciser les relations entre ces deux notions, rappelons d'abord que $\widehat{\mathcal{T}}^{+}=\Sigma$ se déduit de $\widehat{\mathcal{T}}$ comme noyau de $\chi\mid\widehat{\mathcal{T}}:\widehat{\mathcal{T}}\longrightarrow\{\pm1\}$. D'autre part, \deleted{soit} (pour un ss-groupe de Teichmüller géom. strict choisi $\widehat{\mathcal{T}}^{+}$) considérons
\[
(14)\qquad
\unsure{\mathcal{N}|N}_{\Sigma}=\mathcal{N}
=\operatorname{Norm}_{\widehat{\widehat{\mathcal{T}}}}
\bigl(\deleted{\widehat{\mathcal{T}}^{+}}\inserted{\Sigma}\bigr)
\qquad (\Sigma=\widehat{\mathcal{T}}^{+})
\]
(peut-être tjrs égal à $\widehat{\widehat{\mathcal{T}}}$ tt entier)

évidemment (si $\deleted{\widehat{\mathcal{T}}^{+}}\inserted{\Sigma}$ provient d'un $\widehat{\mathcal{T}}$)
\[
(15)\qquad
\underset{\inserted{\unsure{=}\widehat{\mathcal{T}}^{+}}}
{\deleted{\widehat{\mathcal{T}}^{+}}\inserted{\Sigma}}
\subset\widehat{\mathcal{T}}\subset\mathcal{N}_{\Sigma}.
\]
\deleted{et $\mathcal{N}/\widehat{\mathcal{T}}^{+}$ est}
Le groupe profini $\mathcal{N}/\deleted{\widehat{\mathcal{T}}^{+}}\inserted{\Sigma}$
\deleted{est} associé à \deleted{la \unsure{struc}}
$\deleted{\widehat{\mathcal{T}}^{+}}\inserted{\Sigma}$ se note
$\boldsymbol{\Gamma}_{\deleted{\unsure{\widehat{\mathcal{T}}^{+}}}\inserted{\Sigma}}$
\deleted{\illegible{1}} \deleted{\illegible{2}}
(si $\deleted{\widehat{\mathcal{T}}^{+}}\inserted{\Sigma}$ est unique, on note $\boldsymbol{\Gamma}_{\pi}$);

\mspage{323}

\scanpage{325}

et les $\widehat{\mathcal{T}}\deleted{\unsure{\mathcal G}}$ \unsure{donnant} naissance au $\overline{\mathrm m}$ $\deleted{\widehat{\mathcal{T}}^{+}}\inserted{\Sigma}$ correspondent \deleted{\unsure{aux cl}} à une classe de conjugaison d'éléments d'ordre $2$ de $\boldsymbol{\Gamma}_{\deleted{\unsure{\widehat{\mathcal{T}}^{+}}}\inserted{\Sigma}}$. Il est clair \deleted{\unsure{qu'on|que}} en tous cas qu'on a une application injective
\[
(16)\qquad
\begin{gathered}
\text{ens des ss-groupes}\\
\text{de Teichmüller géom}\\
\widehat{\mathcal{T}}\text{ de }\widehat{\widehat{\mathcal{T}}}\text{ tels que}\\
\widehat{\mathcal{T}}^{+}=\Sigma
\end{gathered}
\longrightarrow
\begin{gathered}
\text{ens des }\unsure{\text{éléments}|\text{élémts}}\\
\text{d'ordre }2\text{ de }\boldsymbol{\Gamma}_{\Sigma}
\end{gathered}
\quad :
\]
et que son image est stable par conjugaison. Montrons que \unsure{deux|les} éléments de l'image sont conjugués par $\boldsymbol{\Gamma}_{\Sigma}$. En effet, soient
\[
\widehat{\mathcal{T}},\ \widehat{\mathcal{T}}^{\prime}\supset\Sigma
\quad\text{tels que}\quad
\Sigma=\widehat{\mathcal{T}}^{+}=\widehat{\mathcal{T}}^{\prime+}.
\]
\deleted{On \illegible{1} \unsure{évidemment}}
\inserted{$\exists\,g\in\widehat{\widehat{\mathcal{T}}}$}
tel que
$\widehat{\mathcal{T}}^{\prime}=\operatorname{int}(g)\unsure{\mathcal{T}}$,
et on \unsure{aura|a} donc
$\widehat{\mathcal{T}}^{\prime+}=\operatorname{int}_{g}\widehat{\mathcal{T}}^{+}$
i.e. $g\in\mathcal{N}$, OK.

Les éléments d'ordre $2$ ainsi obtenus dans $\boldsymbol{\Gamma}_{\Sigma}$ s'appellent les involutions canoniques. On en donnera une interprétation conjecturale remarquable plus bas. L'ens. $\widehat{\widehat{\mathcal{T}}}/\mathcal{N}$ s'identifie à l'ens des \unsure{ss-groupes} de

\mspage{324}

\scanpage{326}
\mspage{165}

Teichmüller géométriques stricts (une fois choisi l'élément ``origine'' $\Sigma$). Plus \unsure{intrinsèquement}, on \unsure{aura}:
\[
(17)\qquad
\begin{gathered}
\text{Ens des ss-groupes de}\\
\text{Teichm. géom stricts}\\
\text{de }\widehat{\widehat{\mathcal{T}}}(\pi)
\end{gathered}
\unsure{\xleftarrow{\sim}|\simeq}
\operatorname{Isomext}_{\unsure{\mathrm{lac}}}
(\widehat{\Pi}_{g,\nu},\pi)/\unsure{\mathcal{N}_{g,\nu}}
\]
où on pose, comme de juste
\[
(18)\qquad
\mathcal{N}_{g,\nu}
=\operatorname{Norm}_{\widehat{\widehat{\mathcal{T}}}_{g,\nu}}
\bigl(\widehat{\mathcal{T}}_{g,\nu}^{+}\bigr)
\]
(peut-être égal à $\widehat{\widehat{\mathcal{T}}}_{g,\nu}$ tt entier!)

Considérons maintenant le \uline{topos modulaire} \uline{sur} $\mathbb{Q}$ des courbes algébriques de type $(g,\nu)$, noté $M_{g,\nu\,\mathbb{Q}}$, ou simplement $M_{g,\nu}$. Si nous choisissons un rev. universel $\widetilde{M}_{g,\nu}$ de $M_{g,\nu\,\mathbb{Q}}$ (d'où un rev. universel de $\unsure{\operatorname{Spec}}\mathbb{Q}$ i.e. une clôture alg. $\overline{\mathbb{Q}}$ de $\mathbb{Q}$) on peut \unsure{préciser} le $\pi_1(M_{g,\nu\,\mathbb{Q}})$ comme le groupe des $M_{g,\nu\,\mathbb{Q}}$-\unsure{autom} de ce \unsure{revêt|revêtement}; \deleted{\unsure{et en} \unsure{faire} \illegible{1}} \inserted{l'appeler le} \uline{groupe de Teichmüller arithmétique} de type $(g,\nu)$ (relatif au choix de $\widetilde{M}_{g,\nu\,\mathbb{Q}}$). On aura donc une suite exacte

\mspage{325}

\scanpage{327}

\[
(19)\qquad 1\longrightarrow \pi_1(M_{g,\nu\,\overline{Q}})
\longrightarrow \pi_1(M_{g,\nu})
\longrightarrow \operatorname{Gal}(\overline{Q}/Q)\longrightarrow 1
\]
\deleted{\(1\longrightarrow \unsure{\widehat{\mathcal{T}}_{g,\nu}}\longrightarrow \illegible{1}\)}

(\unsure{où} $\overline{Q}$ et les pts base sont \unsure{explicités} comme il a été dit) Notons que par la théorie \unsure{transcendante} de Teichmüller, on a un iso. (défini mod \unsure{opération} \unsure{induite} par $\pi_1(M_{g,\nu})$ sur $\pi_1(M_{g,\nu,\overline{Q}})$)
\[
(20)\qquad \pi_1(\unsure{M_{g,\nu,\overline{Q}}})\simeq\widehat{\mathcal{T}}_{g,\nu}^{+}
\qquad (\text{compactifié profini de }\mathcal{T}_{g,\nu}^{+})
\]

Considérons d'autre part le \unsure{schéma.} $\unsure{U|\mathcal{U}}_{g,\nu}$\inserted{ sur $M_{g,\nu}$}, courbe de type $(g,\nu)$ ``universelle''

On a donc, en choisissant un rev. universel $\widetilde{\unsure{U|\mathcal{U}}}_{g,\nu}$ de celle-ci \uline{\unsure{en dessus}} du rev. universel choisi $\unsure{\widetilde{M}_{g,\nu}^{\natural}}$ de $M_{g,\nu}$, une suite exacte
\diagram{
\begin{tikzcd}[column sep=large,row sep=large,ampersand replacement=\&]
(21) \& 1 \arrow[r] \& \pi \arrow[r] \arrow[dd,"\Vert" description] \& \pi_1(\unsure{\overset{\deleted{\sim}}{\unsure{U|\mathcal{U}}}_{g,\nu}}) \arrow[r] \arrow[dd] \& \pi_1(M_{g,\nu}) \arrow[r] \arrow[dd] \& 1 \\
\text{qui s'envoie dans :} \& \& \& \& \& \\
(22) \& 1 \arrow[r] \& \pi \arrow[r] \& \widehat{\widehat{\unsure{\mathfrak{S}|\mathcal{G}}}}(\pi) \arrow[r] \& \widehat{\widehat{\mathcal{T}}}(\pi) \arrow[r] \& 1
\end{tikzcd}
}
et de $\overline{\mathrm m}$ pour $\unsure{U|\mathcal{U}}_{g,\nu\,\overline{Q}}\longrightarrow M_{g,\nu\,\unsure{\overline{Q}|Q}}$

\mspage{326}

\scanpage{328}
\mspage{\unsure{166|168}}

\[
\unsure{(\overline{21})|(\overline{21}')}\qquad
1\longrightarrow\pi\longrightarrow\pi_1(\unsure{U|\mathcal{U}}_{g,\nu\,\overline{Q}})
\longrightarrow\pi_1(M_{g,\nu\,\overline{Q}})\longrightarrow1
\]
\deleted{\(\unsure{(\overline{22})}\qquad 1\longrightarrow\pi\longrightarrow\)}

On a bien sur un \unsure{isom} de groupes à \unsure{lacets|bouts}
\[
(23)\qquad \pi\simeq\widehat{\Pi}_{g,\nu}
\]
dont on voudrait déterminer l'indétermination de façon précise

Notons $M_{g,\nu\,C}$ le \deleted{\illegible{1}}\inserted{topos} modulaire de Teichmüller complexe, défini à l'aide de la surface \deleted{diff.}\inserted{$C^\infty$} $\unsure{U|\mathcal{U}}_{g,\nu}$, il est muni du rev. universel canonique $\widetilde{M}_{g,\nu\,C}$ (\S~21)
\margin{$\widetilde{M}_{g,\nu\,C}=E_g/\unsure{A_{g,\nu}^{0}|A_{g,\nu}^{\circ}}$}
de groupe \unsure{d'autom.} $\mathcal{T}_{g,\nu}^{+}$ justement \unsure{--|...}
$\unsure{U|\mathcal{U}}_{g,\nu\,C}$ étant lui aussi muni d'un rev. universel \unsure{en dessus} du précédent,
$\widetilde{\unsure{U|\mathcal{U}}}_{g,\nu\,C}=E_g/\unsure{A_{g,\nu+1}^{0}|A_{g,\nu+1}^{\circ}}$.
Ceci donne naissance à la suite exacte de groupes discrets (dans le contexte topologie générale)
\diagram{
\begin{tikzcd}[column sep=large,row sep=large,ampersand replacement=\&]
(24) \& 1 \arrow[r] \& \Pi_{g,\nu} \arrow[r] \arrow[dd] \& \unsure{\mathfrak{S}|\mathcal{G}}_{g,\nu}^{+} \arrow[r] \arrow[d,equal] \& \mathcal{T}_{g,\nu}^{+} \arrow[r] \arrow[d,equal] \& 1 \\
\& \& \& \pi_1(\unsure{U|\mathcal{U}}_{g,\nu\,C}) \arrow[d] \& \pi_1(M_{g,\nu\,C}) \arrow[d] \& \\
\text{(s'envoyant dans} \& 1 \arrow[r] \& \Pi_{g,\nu} \arrow[r] \& \operatorname{Aut}_{\mathrm{lac}}\deleted{\unsure{\mathrm{ext}}}(\Pi_{g,\nu}) \arrow[r] \arrow[d,equal] \& \operatorname{Aut\,ext}_{\mathrm{lac}}(\Pi_{g,\nu}) \arrow[r] \arrow[d,equal] \& 1 \\
\& \& \& \unsure{\mathfrak{S}|\mathcal{G}}_{g,\nu} \& \mathcal{T}_{g,\nu} \&
\end{tikzcd}
}
par un \unsure{isom|hom.} de \unsure{structures} d'extensions), qui par passage aux complétés profinis donne la suite exacte analogue pour les

\mspage{327}

\scanpage{329}

multiplicités algébriques sur $\mathbb C$

\diagram{\begin{tikzcd}[ampersand replacement=\&,column sep=1.3em,row sep=1.2em]
(25) \& 1 \arrow[r] \&
\widehat{\Pi}_{g,\nu} \arrow[r] \arrow[d,equal] \&
\deleted{\unsure{\widehat{\Pi}}}\,\widehat{\mathcal G}_{g,\nu}
  \arrow[r] \arrow[d,very thick] \&
\widehat{\mathcal T}_{g,\nu}^{+} \arrow[r]
  \arrow[d,phantom,"{\rotatebox{-90}{$\unsure{\hookrightarrow}$}}" description] \& 1 \\
\& \& \inserted{\widehat{\Pi}_{g,\nu}} \arrow[r] \&
\inserted{\operatorname{Aut}\deleted{\unsure{\mathrm{ext}}}_{\mathrm{lac}}(\widehat{\Pi}_{g,\nu})}
  \arrow[r] \&
\inserted{\operatorname{Autext}_{\mathrm{lac}}(\widehat{\Pi}_{g,\nu})}
  \& \\[-0.7em]
\& \& \& |[xshift=-2.2em]| {\inserted{\raisebox{1ex}{\rotatebox[origin=c]{-60}{$=$}}\,\unsure{\widehat{\widehat{\mathcal G}}}_{g,\nu}}} \&
|[xshift=-2.2em]| {\inserted{\raisebox{1ex}{\rotatebox[origin=c]{-60}{$=$}}\,\widehat{\widehat{\mathcal T}}_{g,\nu}}} \&
\end{tikzcd}}

Cette situation \unsure{provient} d'ailleurs \unsure{canoniquement} d'une situation analogue sur $\overline{\mathbb Q}_0=\text{clôture alg de }\mathbb Q\text{ dans }\mathbb C$, le choix de $\widetilde{\unsure{U|\mathcal{U}}}_{g,\nu\mathbb C}$ définissant un $\widetilde{\unsure{U|\mathcal{U}}}_{g,\nu\overline{\mathbb Q}_0}$, d'où un $\widetilde{M}_{g,\nu\overline{\mathbb Q}_0}$.

Pour définir un isom entre cette suite (25) et la suite exacte $(\overline{21})$ \deleted{il faut} (\unsure{munie} de \unsure{$2$ homs|son hom.} dans (22)) il faut

$1^{\circ})$ choisir un isom $\overline{\mathbb Q}\simeq\overline{\mathbb Q}_0$ --- l'indétermination est dans $\mathrm{Gal}(\overline{\mathbb Q}/\mathbb Q)$ ; ce qui permet d'identifier $\overline{\mathbb Q}$ et $\overline{\mathbb Q}_0$ ;

$2^{\circ})$ Choisir un isom (\deleted{\unsure{entre}} $\overline{\mathbb Q}=\overline{\mathbb Q}_0$) de $(\widetilde{\unsure{U|\mathcal{U}}}_{g,\nu})_0$ (construit par voie \unsure{transcendante|transcendant} sur $\mathbb C$ et descendu à $\overline{\mathbb Q}_0$) avec $\widetilde{\unsure{U|\mathcal{U}}}_{g,\nu\overline{\mathbb Q}}$ --- l'indétermination est dans $\pi_1(M_{g,\nu\overline{\mathbb Q}})$.

En résumé, on a une classe d'iso de \unsure{(24)|(21)}\inserted{\unsure{2|?}} et (25), défini mod autom \inserted{\deleted{\illegible{1}} \illegible{1} \unsure{une fois choisi} $\overline{\mathbb Q}\simeq\overline{\mathbb Q}_0$} intérieurs dans \inserted{le groupe \unsure{\uline{profini}}} $\pi_1(\unsure{U|\mathcal{U}}_{g,\nu})$. (\unsure{Ces|Les} isomorphismes \inserted{\unsure{\uline{correspondants}}} transforment la \deleted{structure} \uline{classe de \unsure{discrétification}} \uline{orientée} \unsure{Standard|standard} de $\widehat{\Pi}_{g,\nu}$ en \deleted{\unsure{une classe}} \inserted{exactement} \uline{une} classe de discrétifications orientées de $\pi$.

\mspage{328}

\scanpage{330}
\mspage{167}

\deleted{Ainsi $\pi$ de (21) est \unsure{can.} muni d'une \deleted{\unsure{classe de}} \inserted{\unsure{quasi}} discrétification orientée \inserted{\unsure{$\pi_0^{\natural}$}}, et dans l'hom de $(\overline{21})$ dans (22), on trouve que $\pi_1(\unsure{U|\mathcal{U}}_{g,\nu\overline{\mathbb Q}})$ s'identifie à $\widehat{\unsure{\mathfrak{S}|\mathcal{G}}}(\unsure{\pi_0^{\natural}})$, et $\pi_1(M_{g,\nu})$ à $\widehat{\mathcal T}(\unsure{\pi_0^{\natural}})$ \unsure{$\longrightarrow$}}

Changeant le choix de l'isom en son complexe conjugué, on trouve la \inserted{\unsure{quasi-}} \unsure{discrétification} orientée opposée.

On voit d'autre part que l'image de $\pi_1(M_{g,\nu\overline{\mathbb Q}})\unsure{\subset}\pi_1(M_{g,\nu})$ dans $\widehat{\widehat{\mathcal T}}(\pi)$ n'est autre que le $\widehat{\mathcal T}^{+}$ \unsure{correspondant|correspond.} à une \inserted{quelconque} des classes de quasi-rigidifications discrètes choisies --- le sous-groupe ne dépend pas du choix de cette classe --- ce qui ne fait qu'exprimer que l'image de $\pi_1(M_{g,\nu\overline{\mathbb Q}})$ dans $\widehat{\widehat{\mathcal T}}(\pi)$, \unsure{considéré|considérée} comme sous-groupe de l'image du groupe plus grand $\pi_1(M_{g,\nu})$, y est \unsure{invariante|invariant} --- ce qui résulte du fait que $\pi_1(M_{g,\nu\overline{\mathbb Q}})$ est \unsure{invariant} dans $\pi_1(M_{g,\nu\unsure{\mathbb Q|\overline{\mathbb Q}}})$. Ainsi $\pi$

\mspage{329}
